\documentclass[11pt,a4paper]{article}
\usepackage[T1]{fontenc}
\usepackage[utf8]{inputenc}
\usepackage{amsmath,amsthm}
\usepackage{newtxtext,newtxmath}
\renewcommand{\ttdefault}{lmtt}
\usepackage{microtype}
\usepackage[a4paper,left=29mm,right=29mm,top=27mm,bottom=29mm]{geometry}
\usepackage{booktabs,array}
\usepackage{xcolor}
\usepackage{enumitem}
\usepackage{fancyhdr}
\usepackage{needspace}
\usepackage{hyperref}
\definecolor{inkblue}{RGB}{35,65,86}
\hypersetup{colorlinks=true,linkcolor=inkblue,citecolor=inkblue,urlcolor=inkblue,
  pdftitle={A Goldbach theorem for the Liouville function},pdfauthor={},
  pdfsubject={An elementary proof with a complete Lean formalization},
  pdfkeywords={Liouville function, Goldbach problem, multiplicative functions, Lean}}
\urlstyle{same}
\linespread{1.055}
\setlength{\parindent}{1.15em}
\setlength{\parskip}{2pt}
\setlength{\emergencystretch}{1.5em}
\setlength{\headheight}{14pt}
\pagestyle{fancy}
\fancyhf{}
\fancyhead[L]{\small\itshape A Goldbach theorem for the Liouville function}
\fancyhead[R]{\small\thepage}
\renewcommand{\headrulewidth}{0pt}
\fancypagestyle{firstpage}{\fancyhf{}\fancyfoot[C]{\small\thepage}}
\newtheorem{theorem}{Theorem}[section]
\newtheorem{proposition}[theorem]{Proposition}
\newtheorem{lemma}[theorem]{Lemma}
\newtheorem{corollary}[theorem]{Corollary}
\theoremstyle{remark}
\newtheorem*{remark}{Remark}
\numberwithin{equation}{section}
\newcommand{\Fp}{\mathbb F_p}
\newcommand{\Fpx}{\mathbb F_p^{\times}}
\newcommand{\Z}{\mathbb Z}
\newcommand{\N}{\mathbb N}
\newcommand{\leg}[2]{\left(\frac{#1}{#2}\right)}
\newcommand{\lam}{\lambda}
\newcommand{\code}[1]{\texttt{#1}}
\newcommand{\mainresult}{\ref{thm:main}}

\begin{document}
\thispagestyle{firstpage}
\begin{center}
  {\LARGE\bfseries A Goldbach theorem\par
   \vspace{3pt}for the Liouville function\par}
  \vspace{9pt}
  {\large Two commuting symmetries and a descent}
\end{center}
\vspace{8pt}

\section{The problem and its background}

The Goldbach conjecture asks whether every even integer greater than
two is a sum of two primes. One can weaken the requirement on the summands
while retaining a basic feature of primality: a prime has an odd number of
prime factors.

For a positive integer \(n\), let \(\Omega(n)\) be the number of its prime
factors, counted with multiplicity, and put
\[
  \lam(n)=(-1)^{\Omega(n)}.
\]
This is the \emph{Liouville function}. Thus \(\lam(1)=1\), every prime has
value \(-1\), and
\begin{equation}\label{eq:liouville}
  \lam(ab)=\lam(a)\lam(b),\qquad
  \lam(2n)=\lam(3n)=-\lam(n).
\end{equation}
The condition \(\lam(n)=-1\) also admits many composite integers: for
example, \(\lam(8)=\lam(12)=-1\), so \(20=8+12\) is an admissible sum.

In August 2018, the MathOverflow user Pablo~\cite{MathOverflow} asked
whether every even \(N>2\) has a representation
\begin{equation}\label{eq:question}
  N=a+b,\qquad a,b\ge1,\qquad \lam(a)=\lam(b)=-1.
\end{equation}
Mangerel attributes the question to Shusterman~\cite[Problem~1.1]{MangerelSigns}.
Goldbach's conjecture implies this assertion. The Liouville condition
retains the parity of prime factorization while allowing composite summands.

A closely related question concerns the reflected correlation
\[
  C(N)=\sum_{1\le n<N}\lam(n)\lam(N-n).
\]
The older convolution problem of Corr\'adi and K\'atai asks for the
stronger cancellation statement \(C(N)=o(N)\) as \(N\to\infty\)
\cite{CorradiKatai}; see also the formulation in
\cite[Conjecture~1]{Krishnamoorthy}. Here the issue is whether a specified
sign pair occurs at every prescribed even total.

Mangerel proved, by an elementary argument, that
\begin{equation}\label{eq:nonextreme}
  |C(N)|<N-1\qquad\text{for every }N\ge11,
\end{equation}
and determined that equality in absolute value occurs, for \(N\ge2\),
exactly at \(N\in\{2,3,5,10\}\)~\cite[Theorem~1.2]{MangerelGoldbach}.
This ensures the occurrence of both equal and unequal signs, but does not
specify which equal-sign pattern occurs. In particular, it does not by
itself give the pair \((-1,-1)\) in~\eqref{eq:question}.

Mangerel subsequently proved~\eqref{eq:question} for every sufficiently
large even integer under the Generalised Riemann Hypothesis for Dirichlet
\(L\)-functions~\cite[Theorem~1.2]{MangerelSigns}. His argument also gives,
for sufficiently large prime totals, a quantitative lower bound for each
of the four reflected sign patterns. A preliminary manuscript treated
positive multiples of four unconditionally, using
\eqref{eq:nonextreme} together with elementary sign
propagation~\cite{MultiplesFour}.

There are also results averaged over the total. For example,
Krishnamoorthy~\cite[Theorem~2]{Krishnamoorthy} gives, for each fixed
\(\delta>0\), a density-zero exceptional set outside which
\(|C(N)|\le\delta N\), by classical exponential-sum and mean-square
methods. Such averaged control allows exceptional totals; the theorem
below concerns every even total. It does not establish the stronger
asymptotic formula \(C(N)=o(N)\).

The result proved here has no size threshold or analytic hypothesis.

\begin{theorem}\label{thm:main}
For every even integer \(N>2\), there are positive integers \(a,b\) such that
\[
  N=a+b\qquad\text{and}\qquad \lam(a)=\lam(b)=-1.
\]
\end{theorem}

The main work goes into an assertion with the opposite sign.

\begin{theorem}\label{thm:prime}
For every prime \(p>3\), there are positive integers \(u,v\) such that
\[
  2p=u+v\qquad\text{and}\qquad \lam(u)=\lam(v)=1.
\]
\end{theorem}

The usefulness of positive-positive pairs at twice a prime was already
pointed out by Thomas Bloom in the original discussion~\cite{MathOverflow}.
Multiplication by an integer of negative Liouville sign turns such a pair
into a negative-negative pair. Section~\ref{sec:finish} gives the short
reduction, including the small prime factors.

\paragraph{The idea of the proof.}
Suppose a positive-positive pair is missing at \(2p\). This forces
one-sided reflection inequalities for \(\lam\). Extend its values below
\(p/2\) oddly to \(\Fp\). The inequalities almost make multiplication by
\(-2\) and \(-3\) preserve this extension: the failures can occur only in
two short intervals. Since the two multiplications commute, these failures
must satisfy a compatibility identity. The intervals leave no room for a
nonzero failure.

The completed function now has exact multiplication laws for \(-1,2,3\),
and it is already multiplicative whenever a positive integer product is
less than \(p/2\). A descent shows that these facts force full
multiplicativity. Finally, a small prime which is a quadratic residue
contradicts agreement with \(\lam\).

The strategy of deriving character-like rigidity from missing sign
patterns is already present in Mangerel's
work~\cite{MangerelGoldbach,MangerelSigns}. Commuting dilation defects also
appear explicitly in~\cite[Section~2.2.2 and Lemma~4.3]{MangerelSigns}.
Here the one-sided missing-pair condition yields exact symmetries through
their interval supports, and a half-interval descent completes the
argument. The proof below needs neither correlation estimates nor
classification of multiplicative characters. Its only classical
number-theoretic input beyond elementary arithmetic is quadratic
reciprocity and its supplementary laws.

\section{Reduction to twice a prime}\label{sec:finish}

We first deduce Theorem~\ref{thm:main} from Theorem~\ref{thm:prime}.
This explains the change of sign in the auxiliary problem. Only two prime
factors need to be extracted.

\begin{proof}[Proof of Theorem~\ref{thm:main}, assuming Theorem~\ref{thm:prime}]
Write \(N=2m>2\). If \(\lam(m)=-1\), use \(N=m+m\).
Otherwise \(\lam(m)=1\). Since \(m>1\), it has a positive even number
of prime factors, counted with multiplicity. Extract two of them:
\[
  m=rsc,\qquad r,s\text{ prime},\qquad c\ge1,\qquad\lam(c)=1.
\]
The primes \(r,s\) need not be distinct.

If one of them, say \(r\), exceeds \(3\), Theorem~\ref{thm:prime}
gives \(2r=u+v\) with \(\lam(u)=\lam(v)=1\). Multiply by \(sc\).
Since \(\lam(sc)=-1\),
\[
  N=scu+scv,\qquad \lam(scu)=\lam(scv)=-1.
\]
Otherwise \(r,s\in\{2,3\}\), and \(2rs\in\{8,12,18\}\).
Use the corresponding identity
\[
  8=3+5,\qquad 12=5+7,\qquad 18=7+11,
\]
and multiply both summands by \(c\). Its positive Liouville sign
preserves their negative signs.
\end{proof}

\section{A missing sum and two commuting symmetries}\label{sec:symmetry}

Fix a prime \(p>3\). Throughout this section, suppose that \(2p\) has no
positive-positive Liouville representation, and abbreviate \(f(n)=\lam(n)\).
Only the values with \(0<n<2p\) will be needed.

\subsection{Reflection becomes an equality on the upper third}

There can be no negative-negative pair summing to \(p\): doubling both
summands would give a positive-positive pair summing to \(2p\). This simple
observation gives all the inequalities needed below.

\begin{lemma}\label{lem:reflection}
For \(0<n<p\),
\begin{equation}\label{eq:reflection}
  f(p-n)\ge-f(n),\qquad f(p+n)\le f(n).
\end{equation}
For \(0<n<p/2\),
\begin{equation}\label{eq:double-reflection}
  f(p-2n)\ge f(n).
\end{equation}
\end{lemma}

\begin{proof}
The first inequality in~\eqref{eq:reflection} says precisely that \(n\)
and \(p-n\) cannot both have negative sign. For the second, the only
nontrivial case is \(f(n)=-1\). The first inequality then gives
\(f(p-n)=1\); since \((p-n)+(p+n)=2p\), the missing positive-positive
pair forces \(f(p+n)=-1\).
Finally, apply the first inequality at \(2n\) and use
\(f(2n)=-f(n)\).
\end{proof}

The inequalities are ordered as ordinary integers, with \(-1<1\). Their
useful feature is that a closed chain forces equality.

\begin{lemma}\label{lem:upper}
If \(p/3<n<p/2\), then
\begin{equation}\label{eq:upper}
  f(n)=f(p-2n)=-f(3n-p).
\end{equation}
\end{lemma}

\begin{proof}
Put \(u=p-2n\) and \(v=3n-p\); both are positive. By
\eqref{eq:double-reflection}, \(f(n)\le f(u)\). Since \(3n=p+v\),
translation in~\eqref{eq:reflection} gives \(-f(n)\le f(v)\).
Finally, \(p=3u+2v\). The two terms \(3u,2v\) cannot both have
negative sign, so \(f(u),f(v)\) cannot both have positive sign.
Equivalently, \(f(v)\le-f(u)\). Thus
\[
  -f(n)\le f(v)\le-f(u)\le-f(n),
\]
and every inequality is an equality.
\end{proof}

\subsection{Where can the symmetries fail?}

Every nonzero residue \(x\in\Fp\) has a unique integer representative
\(t\) with \(-p/2<t<p/2\). Define
\begin{equation}\label{eq:completion}
  G(0)=0,\qquad G(x)=\operatorname{sgn}(t)f(|t|)\quad(x\ne0).
\end{equation}
Thus \(G\) is odd, \(G(1)=1\), and \(G(n)=f(n)\) for
\(0<n<p/2\). Integer arguments of \(G\) always mean their residues
modulo \(p\). We call \(t\) the \emph{central representative} of \(x\).

If \(0<n<p/4\), local doubling already gives
\(G(-2n)=-f(2n)=f(n)=G(n)\). Lemma~\ref{lem:upper} gives the same
identity for \(p/3<n<p/2\). Only the interval between \(p/4\) and
\(p/3\) remains. To treat doubling and tripling together, define their
\emph{defects}
\begin{equation}\label{eq:defects}
  A(x)=G(-2x)-G(x),\qquad B(x)=G(-3x)-G(x).
\end{equation}
Both functions are odd. They record exactly what still needs to be proved.

\Needspace{10\baselineskip}
\begin{lemma}\label{lem:support}
For positive central representatives, the defects satisfy the following
table. An interval in the table means its integer points.
\begin{center}
\renewcommand{\arraystretch}{1.3}
\begin{tabular}{@{}lcccc@{}}
\toprule
\(n\) & \((0,p/6)\) & \((p/6,p/4)\) & \((p/4,p/3)\) & \((p/3,p/2)\) \\
\midrule
\(A(n)\) & \(0\) & \(0\) & \(\ge0\) & \(0\) \\
\(B(n)\) & \(0\) & \(\ge0\) & \(A(n)\) & \(0\) \\
\bottomrule
\end{tabular}
\end{center}
\end{lemma}

\begin{proof}
Local doubling and Lemma~\ref{lem:upper} give the zeros in the first
row. On the remaining interval,
\[
  A(n)=f(p-2n)-f(n)\ge0
\]
by~\eqref{eq:double-reflection}. Local tripling and
Lemma~\ref{lem:upper} give the two outer zeros in the second row.
For \(p/6<n<p/3\), the central representative of \(-3n\) is
\(p-3n\), and reflection at \(3n\) gives
\[
  B(n)=f(p-3n)-f(n)\ge0.
\]
To compare the two rows when \(p/4<n<p/3\), apply
Lemma~\ref{lem:upper} to \(u=p-2n\in(p/3,p/2)\):
\[
  f(p-2n)=-f\bigl(2(p-3n)\bigr)=f(p-3n).
\]
Subtracting \(f(n)\) gives \(A(n)=B(n)\). There are no missing
boundary cases, since none of \(p/6,p/4,p/3,p/2\) is an integer.
\end{proof}

\subsection{The commuting square}

The table leaves one extra interval on which \(B\) may be nonzero.
Commutativity removes it.

\begin{proposition}\label{prop:symmetries}
The completion \(G\) satisfies, for every \(x\in\Fp\),
\begin{equation}\label{eq:exact}
  G(2x)=-G(x),\qquad G(3x)=-G(x).
\end{equation}
\end{proposition}

\begin{proof}
Multiplication by \(-2\) commutes with multiplication by \(-3\).
Taking the difference of \(G\) along the two routes from \(x\) to
\(6x\) gives
\begin{equation}\label{eq:commuting}
  A(-3x)+B(x)=B(-2x)+A(x).
\end{equation}
Indeed, both sides telescope to \(G(6x)-G(x)\).

Take \(p/6<n<p/4\). The right side of~\eqref{eq:commuting} is zero:
\(A(n)=0\), while \(-2n\) has central magnitude
\(2n\in(p/3,p/2)\), where \(B\) vanishes. On the left,
\(-3n\) has positive central representative
\(p-3n\in(p/4,p/2)\). Hence \(A(-3n)\ge0\) and \(B(n)\ge0\).
Their sum is zero, so \(B(n)=0\).

The two rows of the table now agree on every positive representative.
Oddness gives \(A=B\) on all of \(\Fp\). By~\eqref{eq:defects},
\(G(2x)=G(3x)\), or equivalently
\begin{equation}\label{eq:ratio}
  G\bigl((2/3)x\bigr)=G(x).
\end{equation}
Here \(2/3\) denotes \(2\cdot3^{-1}\) in \(\Fp\). Applying
\eqref{eq:ratio} at \(x\) and at \(-2x\) shows that \(A\), too,
is invariant under multiplication by \(2/3\).

Suppose \(A\) is not identically zero. The table and oddness give a
positive central representative \(n\in(p/4,p/3)\) with \(A(n)>0\).
Let \(m\) be the central representative of \((2/3)n\). Invariance
gives \(A(m)=A(n)>0\). It follows from oddness and the table that
\(m\) is also positive and lies in \((p/4,p/3)\). Yet
\[
  3m\equiv2n\pmod p,
\]
whereas
\[
  2n\in(p/2,2p/3),\qquad 3m\in(3p/4,p).
\]
These are disjoint intervals inside \((0,p)\), so the congruence is
impossible. Therefore \(A=B=0\). The desired identities follow from
\eqref{eq:defects} and oddness of \(G\).
\end{proof}

\section{From short products to full multiplicativity}\label{sec:extension}

The completed function already knows how to multiply whenever reduction
modulo \(p\) is unnecessary. In fact,
\begin{equation}\label{eq:short}
  G(ab)=G(a)G(b)\qquad(a,b\ge1,\ ab<p/2),
\end{equation}
because \(a,b,ab\) are all in the interval where \(G=\lam\).
We now show that exact multiplication by \(-1,2,3\) is enough to carry
this local rule through every reduction modulo \(p\).

The arithmetic step is most transparent as a statement about cosets.
For a subgroup \(H\subseteq\Fpx\), a coset \(xH\) is the set obtained
by multiplying \(x\) by every element of \(H\).
Throughout this section, integers appearing as arguments of a function on
\(\Fpx\), or in a statement of subgroup membership, are reduced modulo \(p\).

\begin{lemma}[Short representatives]\label{lem:cosets}
Let \(p>3\) be prime, and let \(H\) be a proper subgroup of
\(\Fpx\) containing \(-1,2,3\). Let \(q\) be the least positive
integer representing a nonzero residue outside \(H\). Every coset of \(H\)
has a positive integer representative \(n\) with
\[
  n<\frac{p}{2q}.
\]
\end{lemma}

\begin{proof}
First, \(q\) is prime: if it factored into two smaller positive
integers, both factors would belong to \(H\). Since \(1,2,3,4\in H\),
we have \(q\ge5\). Since \(-1\in H\), a missing residue has, up to
sign, a missing positive representative below \(p/2\); thus \(q<p/2\).

Choose the smallest positive integer representative \(n\) of a fixed
coset. Negation gives \(n<p/2\), and \(n\) is odd: an even \(n\)
could be replaced by \(n/2\), because \(2\in H\). If \(n=1\), the
claim follows from \(q<p/2\). Hence assume \(1<n<p/2\), so
\(n\nmid p\). Suppose for a contradiction that \(2qn\ge p\).
Equality is impossible because \(p\) is odd, so \(2qn>p\).

If \(p<(2q-1)n\), choose an even integer \(a\) nearest to \(p/n\).
Then
\[
  2\le a\le2q-2,\qquad 0<|p-an|<n.
\]
The strict inequalities follow because \(p/n\) is not an integer;
in particular, it cannot lie halfway between two consecutive even
integers. Write \(a=2j\), where \(1\le j<q\). Both \(2\) and
\(j\) belong to \(H\), so \(a\in H\). Therefore \(|p-an|\) is a
smaller positive representative of the same coset, a contradiction.

The endpoint \(p=(2q-1)n\) is excluded by \(n\nmid p\). Only the narrow interval
\begin{equation}\label{eq:strip}
  (2q-1)n<p<2qn
\end{equation}
remains. Since the prime \(q\ge5\) is not divisible by \(3\), exactly
one of \(2q-1\) and \(2q+1\) is divisible by \(3\). Choose that
integer as \(a\). It has the form \(a=3j\) with \(1\le j<q\),
so once again \(a\in H\). The bounds~\eqref{eq:strip} give
\[
  0<|p-an|<2n.
\]
This remainder might be larger than \(n\), but it is even: \(a,n,p\)
are all odd. Consequently
\begin{equation}\label{eq:halved}
  r=\frac{|p-an|}{2}
\end{equation}
is a positive integer smaller than \(n\). Its residue is
\(\pm(a/2)n\), which lies in the same coset because
\(-1,a,2\in H\). This final contradiction proves the lemma.
\end{proof}

The last step is the reason for using both small multipliers. Multiples
of \(2\) handle all but the interval~\eqref{eq:strip}; a nearby multiple
of \(3\), followed by division by \(2\), handles the interval that remains.

\begin{proposition}[Half-interval extension]\label{prop:extension}
Let \(p>3\) be prime, and let \(F:\Fpx\to\{\pm1\}\) satisfy
\(F(1)=1\). Suppose that
\begin{align}
  F(cx)&=F(c)F(x)&&\text{for }c\in\{-1,2,3\},\ x\in\Fpx,
    \label{eq:small-exact}\\
  F(ab)&=F(a)F(b)&&\text{for }a,b\in\Z_{\ge1},\ ab<p/2.
    \label{eq:short-F}
\end{align}
Then \(F(xy)=F(x)F(y)\) for every \(x,y\in\Fpx\).
\end{proposition}

\begin{proof}
Call a residue \(c\in\Fpx\) \emph{exact} if \(F(cx)=F(c)F(x)\) for every
\(x\in\Fpx\), and let \(H\) be the set of exact residues.
It contains \(1\) and is closed under multiplication, by composing the
two multiplication laws. It is also closed under inverses: exactness at
\(c\), applied first to \(c^{-1}\) and then to \(c^{-1}x\), gives
\(F(c^{-1}x)=F(c^{-1})F(x)\). Thus \(H\) is a subgroup containing
\(-1,2,3\).

Suppose \(H\) is proper, and let \(q\) be its least missing positive
integer. The multiplicative defect
\[
  D(x)=\frac{F(qx)}{F(q)F(x)}
\]
is constant on each coset of \(H\). Indeed, for \(c\in H\), exactness
at \(c\) cancels the same factor \(F(c)\) from numerator and denominator,
giving \(D(cx)=D(x)\).

By Lemma~\ref{lem:cosets}, every coset contains an integer
\(n<p/(2q)\). Then \(qn<p/2\), so~\eqref{eq:short-F} gives
\(D(n)=1\). Hence \(D=1\) everywhere, which means that \(q\) is
exact after all. Therefore \(H=\Fpx\).
\end{proof}

\begin{corollary}\label{cor:multiplicative}
Under the missing-pair assumption of Section~\ref{sec:symmetry}, the
completion \(G\) is multiplicative on \(\Fpx\).
\end{corollary}

\begin{proof}
Oddness gives the exact multiplier \(-1\). Evaluating
\eqref{eq:exact} at \(x=1\) gives \(G(2)=G(3)=-1\), so the same
identities give exact multiplication by \(2\) and \(3\).
Equation~\eqref{eq:short} supplies short multiplicativity. Apply
Proposition~\ref{prop:extension}.
\end{proof}

\section{A small prime quadratic residue}\label{sec:residue}

Every multiplicative sign function takes the value \(1\) on nonzero
squares. A single prime which is a square modulo \(p\) will therefore
contradict agreement with \(\lam\). The half interval is large enough
to contain such a prime whenever the sign at \(-1\) is negative.

We use the usual quadratic-residue criteria and quadratic reciprocity;
see, for example, Ireland and Rosen~\cite[Chapter~5]{IrelandRosen}.
The prime-divisor observation in the next lemma is standard; it also
appears in the official HMMT 2020 solutions~\cite[Problem~9]{HMMT}.
For the classical study of the least prime quadratic residue, see
Chowla, Cowles and Cowles~\cite{ChowlaCowlesCowles}.

\begin{lemma}\label{lem:small-prime}
If \(p>3\) is prime and \(p\equiv3\pmod4\), then there is a prime
\(\ell<p/2\) which is a quadratic residue modulo \(p\).
In fact, every prime divisor of \((p+1)/4\) has this property.
\end{lemma}

\begin{proof}
The integer \(L=(p+1)/4\) is greater than one. Choose a prime
\(\ell\mid L\); then \(\ell\le L<p/2\). If \(\ell=2\), then
\(L\) is even and \(p\equiv7\pmod8\), so \(2\) is a quadratic
residue modulo \(p\).

If \(\ell\) is odd, quadratic reciprocity and \(p\equiv3\pmod4\)
give
\[
  \leg{\ell}{p}
  =(-1)^{(\ell-1)/2}\leg{p}{\ell}
  =\leg{-p}{\ell}
  =\leg{1}{\ell}=1,
\]
since \(\ell\mid p+1\). Here the parentheses denote the Legendre
symbol. Thus \(\ell\) is a nonzero square modulo \(p\) in either case.
\end{proof}

\begin{proof}[Proof of Theorem~\ref{thm:prime}]
Suppose the asserted pair is missing. The preceding sections construct an
odd multiplicative function \(G:\Fpx\to\{\pm1\}\) agreeing with
\(\lam\) on \(0<n<p/2\).

For every nonzero \(x\),
\[
  G(x^2)=G(x)^2=1.
\]
But \(G(-1)=-1\), so \(-1\) is not a square modulo \(p\).
The standard criterion gives \(p\equiv3\pmod4\).
Choose the prime \(\ell\) from Lemma~\ref{lem:small-prime}. Since
\(\ell\) is a nonzero square and \(\ell<p/2\),
\[
  1=G(\ell)=\lam(\ell)=-1,
\]
a contradiction.
\end{proof}

\clearpage
\appendix
\section{The accompanying formal proof}\label{sec:lean}

Theorem~\ref{thm:main} is formalized in Lean~4~\cite{Lean4} with
mathlib~\cite{Mathlib}. The final declaration is
\begin{center}
\small\code{LiouvilleGoldbach.liouville\_goldbach}.
\end{center}
It has exactly the stated hypotheses and conclusion, and uses mathlib's
Liouville function on natural numbers.

The principal steps are in \code{Commuting.lean}, \code{Descent.lean},
\code{Coset.lean}, and \code{Residue.lean}. The third module proves the
short-representative lemma in a more general invariant-set form and
deduces the same extension theorem. \code{Reduction.lean} and
\code{Main.lean} complete the argument.

The source archive pins Lean to \code{v4.34.0-rc2} and mathlib to commit
\begin{center}
\small\code{de2ef68216c6074f338c8e61890ee0a379ddfb9b}.
\end{center}
The complete build passes. The final axiom audit reports only
\code{propext}, \code{Classical.choice}, and \code{Quot.sound}, the
standard logical foundations used by Lean and mathlib. No unfinished
proofs or additional mathematical axioms occur. A dependency audit also
checks that the final theorem uses the new commuting and descent
arguments. The accompanying archive provides the source, reproduction
instructions, and all verification logs.

\begin{thebibliography}{11}
\small
\setlength{\itemsep}{4pt}

\bibitem{MathOverflow}
Pablo (MathOverflow username),
\emph{Goldbach's conjecture for the Liouville function},
MathOverflow discussion, question 307479, 3 August 2018.
\href{https://mathoverflow.net/questions/307479/goldbachs-conjecture-for-the-liouville-function}{mathoverflow.net/questions/307479}.

\bibitem{CorradiKatai}
C.~A.~Corr\'adi and I.~K\'atai,
\emph{Some problems concerning the convolutions of number-theoretical
functions}, Archiv der Mathematik \textbf{20} (1969), 24--29.
\href{https://doi.org/10.1007/BF01898987}{doi:10.1007/BF01898987}.

\bibitem{MangerelGoldbach}
A.~P.~Mangerel,
\emph{On a Goldbach-type problem for the Liouville function},
International Mathematics Research Notices \textbf{2024} (2024), no.~16,
11865--11877.
\href{https://doi.org/10.1093/imrn/rnae149}{doi:10.1093/imrn/rnae149}.

\bibitem{MangerelSigns}
A.~P.~Mangerel,
\emph{On Shusterman's Goldbach-type problem for sign patterns of the
Liouville function}, preprint, 23 December 2024.
\href{https://arxiv.org/abs/2412.17199v1}{arXiv:2412.17199v1}.

\bibitem{Krishnamoorthy}
K.~Krishnamoorthy,
\emph{On a conjecture of Corr\'adi and K\'atai}, preprint, 2026.
\href{https://arxiv.org/abs/2608.13266}{arXiv:2608.13266}.

\bibitem{MultiplesFour}
\emph{Liouville--Goldbach for multiples of four}, preliminary manuscript
and formalization, CaptainSude GitHub repository, accessed 16 September 2026.
\href{https://github.com/CaptainSude/Liouville-Goldbach-Multiples-of-Four}{Repository and manuscript}.

\bibitem{IrelandRosen}
K.~Ireland and M.~Rosen,
\emph{A Classical Introduction to Modern Number Theory}, 2nd ed.,
Graduate Texts in Mathematics, vol.~84, Springer, New York, 1990.
\href{https://doi.org/10.1007/978-1-4757-2103-4}{doi:10.1007/978-1-4757-2103-4}.

\bibitem{HMMT}
Harvard--MIT Mathematics Tournament,
\emph{HMMT February 2020: Team Round Solutions}, 15 February 2020,
Problem~9, solution~1, p.~8; problem proposed by A.~Gu.
\href{https://hmmt-archive.s3.amazonaws.com/tournaments/2020/feb/team/solutions.pdf}{Official solutions}.

\bibitem{ChowlaCowlesCowles}
S.~Chowla, J.~Cowles and M.~Cowles,
\emph{The least prime quadratic residue and the class number},
Journal of Number Theory \textbf{22} (1986), no.~1, 1--3.
\href{https://doi.org/10.1016/0022-314X(86)90026-0}{doi:10.1016/0022-314X(86)90026-0}.

\bibitem{Lean4}
L.~de~Moura and S.~Ullrich,
\emph{The Lean 4 theorem prover and programming language},
in Automated Deduction---CADE~28, Lecture Notes in Computer Science,
vol.~12699, Springer, 2021, 625--635.
\href{https://doi.org/10.1007/978-3-030-79876-5_37}{doi:10.1007/978-3-030-79876-5\_37}.

\bibitem{Mathlib}
The mathlib Community,
\emph{The Lean mathematical library},
Proceedings of the 9th ACM SIGPLAN International Conference on Certified
Programs and Proofs (CPP~2020), ACM, 2020, 367--381.
\href{https://doi.org/10.1145/3372885.3373824}{doi:10.1145/3372885.3373824}.

\end{thebibliography}
\end{document}
